\section{Key Management dan Lifecycle}

\subsection{Fundamental Key Management}

Key management adalah proses lengkap untuk mengelola cryptographic keys sepanjang lifecycle mereka \cite{nist800_57}. Key management yang buruk dapat mengkompromikan cipher yang paling kuat sekalipun.

\textbf{Key Lifecycle Stages}:
\begin{enumerate}
  \item \textbf{Generation}: Membuat kunci yang aman dan acak
  \item \textbf{Distribution}: Mengirim kunci ke pihak yang berwenang
  \item \textbf{Storage}: Menyimpan kunci dengan aman
  \item \textbf{Usage}: Menggunakan kunci sesuai kebijakan
  \item \textbf{Rotation}: Mengganti kunci secara berkala
  \item \textbf{Revocation}: Mencabut kunci yang dikompromikan
  \item \textbf{Destruction}: Menghapus kunci secara permanen
\end{enumerate}

\subsection{Cryptographically Secure Random Number Generation}

Kunci kriptografi harus dihasilkan menggunakan CSPRNG (Cryptographically Secure Pseudo-Random Number Generator) \cite{nist800_90a}.

\subsubsection{NIST SP 800-90A Standards}

\begin{lstlisting}[style=cstyle, caption={CSPRNG Implementation}]
#include <stdint.h>
#include <openssl/rand.h>

// NIST SP 800-90A Hash-based DRBG
typedef struct {
    uint8_t value[64];     // Internal state
    uint8_t reseed_counter; // Reseed counter
    uint8_t security_strength; // 128, 192, or 256 bits
} HashDRBG;

// Initialize DRBG with entropy
void hash_drbg_init(HashDRBG *drbg, const uint8_t *entropy, size_t entropy_len) {
    // Concatenate entropy, nonce, and personalization string
    uint8_t seed_material[64];
    memcpy(seed_material, entropy, entropy_len);
    
    // Add nonce (timestamp or counter)
    uint64_t nonce = get_timestamp();
    memcpy(seed_material + entropy_len, &nonce, 8);
    
    // Hash to create initial state
    sha256(seed_material, entropy_len + 8, drbg->value);
    
    drbg->reseed_counter = 1;
    drbg->security_strength = 256; // 256-bit security
}

// Generate random bytes
int hash_drbg_generate(HashDRBG *drbg, uint8_t *output, size_t output_len,
                      const uint8_t *additional_input, size_t add_input_len) {
    // Check reseed requirement
    if (drbg->reseed_counter > 10000) {
        return -1; // Need reseed
    }
    
    // Generate pseudo-random bits
    uint8_t temp[64];
    
    // Hash-based generation function
    sha256(drbg->value, 64, temp);
    
    // Copy output
    memcpy(output, temp, output_len);
    
    // Update internal state
    uint8_t update_input[64 + add_input_len];
    memcpy(update_input, temp, 64);
    if (additional_input) {
        memcpy(update_input + 64, additional_input, add_input_len);
    }
    
    sha256(update_input, 64 + add_input_len, drbg->value);
    
    drbg->reseed_counter++;
    return 0;
}

// Reseed with fresh entropy
void hash_drbg_reseed(HashDRBG *drbg, const uint8_t *entropy, size_t entropy_len) {
    uint8_t seed_material[64 + 64]; // Current state + new entropy
    
    memcpy(seed_material, drbg->value, 64);
    memcpy(seed_material + 64, entropy, entropy_len);
    
    sha256(seed_material, 64 + entropy_len, drbg->value);
    drbg->reseed_counter = 1;
}

// Usage example
void generate_aes_key() {
    HashDRBG drbg;
    uint8_t entropy[32];
    uint8_t aes_key[32];
    
    // Get entropy from system
    get_entropy(entropy, 32);
    
    // Initialize DRBG
    hash_drbg_init(&drbg, entropy, 32);
    
    // Generate AES-256 key
    if (hash_drbg_generate(&drbg, aes_key, 32, NULL, 0) == 0) {
        printf("Generated AES-256 key: ");
        for (int i = 0; i < 32; i++) {
            printf("%02x", aes_key[i]);
        }
        printf("\n");
    }
}
\end{lstlisting}

\subsubsection{Entropy Sources}

\begin{table}[htbp]
  \centering
  \small
  \begin{tabular}{ll}
    \toprule
    Entropy Source & Quality \\
    \midrule
    Hardware RNG (Intel RDRAND) & Tinggi \\
    /dev/urandom (Linux) & Tinggi \\
    CryptGenRandom (Windows) & Tinggi \\
    Mouse/keyboard timing & Sedang \\
    Network timing & Rendah \\
    \bottomrule
  \end{tabular}
  \caption{Kualitas entropy sources.}
\end{table}

\subsection{Key Derivation Functions (KDF)}

KDF mengubah kunci master atau password menjadi kunci-kunci turunan dengan properti yang diinginkan \cite{nist800_108}.

\subsubsection{HKDF (HMAC-based Key Derivation Function)}

\begin{lstlisting}[style=cstyle, caption={HKDF Implementation}]
#include <openssl/hmac.h>

// HKDF Extract phase
void hkdf_extract(const uint8_t *salt, size_t salt_len,
                  const uint8_t *ikm, size_t ikm_len,
                  uint8_t *prk) {
    // If salt is empty, use zeros
    uint8_t actual_salt[salt_len ? salt_len : 32];
    if (salt_len == 0) {
        memset(actual_salt, 0, 32);
    } else {
        memcpy(actual_salt, salt, salt_len);
    }
    
    // PRK = HMAC(salt, IKM)
    HMAC(EVP_sha256(), actual_salt, salt_len ? salt_len : 32,
         ikm, ikm_len, prk, NULL);
}

// HKDF Expand phase
void hkdf_expand(const uint8_t *prk, size_t prk_len,
                 const uint8_t *info, size_t info_len,
                 uint8_t *okm, size_t L) {
    uint8_t t[32];  // HMAC output
    uint8_t t_prev[32] = {0};
    int iterations = (L + 31) / 32;
    
    for (int i = 1; i <= iterations; i++) {
        // T(i) = HMAC(PRK, T(i-1) || info || i)
        uint8_t input[32 + info_len + 1];
        int input_len = 0;
        
        if (i > 1) {
            memcpy(input, t_prev, 32);
            input_len += 32;
        }
        
        memcpy(input + input_len, info, info_len);
        input_len += info_len;
        
        input[input_len] = i;
        input_len++;
        
        HMAC(EVP_sha256(), prk, prk_len, input, input_len, t, NULL);
        
        // Copy to output
        int copy_len = (L - (i-1)*32 < 32) ? (L - (i-1)*32) : 32;
        memcpy(okm + (i-1)*32, t, copy_len);
        
        memcpy(t_prev, t, 32);
    }
}

// Complete HKDF
void hkdf(const uint8_t *salt, size_t salt_len,
          const uint8_t *ikm, size_t ikm_len,
          const uint8_t *info, size_t info_len,
          uint8_t *okm, size_t L) {
    uint8_t prk[32];
    
    hkdf_extract(salt, salt_len, ikm, ikm_len, prk);
    hkdf_expand(prk, 32, info, info_len, okm, L);
}

// Usage example: Derive multiple keys from master key
void derive_keys_from_master() {
    uint8_t master_key[32];
    uint8_t salt[16];
    uint8_t derived_keys[64]; // Two 32-bit keys
    
    // Get master key and salt
    get_entropy(master_key, 32);
    get_entropy(salt, 16);
    
    // Derive encryption key and MAC key
    hkdf(salt, 16, master_key, 32, 
         (uint8_t*)"encryption key", 15, derived_keys, 32);
    hkdf(salt, 16, master_key, 32, 
         (uint8_t*)"mac key", 7, derived_keys + 32, 32);
    
    printf("Encryption key: ");
    for (int i = 0; i < 32; i++) printf("%02x", derived_keys[i]);
    printf("\nMAC key: ");
    for (int i = 32; i < 64; i++) printf("%02x", derived_keys[i]);
    printf("\n");
}
\end{lstlisting}

\subsection{Key Storage and Protection}

\subsubsection{Hardware Security Modules (HSM)}

HSM adalah perangkat keras yang dirancang khusus untuk melindungi dan mengelola kunci kriptografi \cite{nist800_57}.

\begin{lstlisting}[style=cstyle, caption={HSM Interface}]
typedef struct {
    int hsm_id;
    uint8_t master_key[32];
    bool tamper_evident;
} HSM;

// HSM key generation (inside secure boundary)
int hsm_generate_key(HSM *hsm, uint8_t *key_id, uint8_t *public_key) {
    if (!hsm->tamper_evident) {
        return -1; // HSM compromised
    }
    
    // Generate key inside HSM
    uint8_t private_key[32];
    generate_key_pair(private_key, public_key);
    
    // Store private key encrypted with master key
    uint8_t encrypted_private[32];
    aes_encrypt(private_key, 32, hsm->master_key, encrypted_private);
    
    // Store in secure memory
    secure_store_key(key_id, encrypted_private, 32);
    
    return 0;
}

// HSM signing operation
int hsm_sign(HSM *hsm, uint8_t *key_id, uint8_t *message, size_t msg_len,
              uint8_t *signature) {
    if (!hsm->tamper_evident) {
        return -1;
    }
    
    // Retrieve and decrypt private key
    uint8_t encrypted_private[32];
    uint8_t private_key[32];
    
    retrieve_key(key_id, encrypted_private, 32);
    aes_decrypt(encrypted_private, 32, hsm->master_key, private_key);
    
    // Perform signing inside HSM
    ecdsa_sign(private_key, message, msg_len, signature);
    
    // Clear sensitive data
    secure_zero_memory(private_key, 32);
    
    return 0;
}
\end{lstlisting}

\subsubsection{Software Key Protection}

\begin{lstlisting}[style=cstyle, caption={Software Key Protection}]
#include <windows.h>
#include <dpapi.h>

// Windows DPAPI for key protection
void protect_key_dpapi(const uint8_t *key, size_t key_len, uint8_t *protected) {
    DATA_BLOB blob_in, blob_out;
    
    blob_in.pbData = (BYTE*)key;
    blob_in.cbData = key_len;
    
    // Encrypt with user-specific key
    if (CryptProtectData(&blob_in, L"Encryption Key", NULL, NULL, NULL, 
                          CRYPTPROTECT_UI_FORBIDDEN, &blob_out)) {
        memcpy(protected, blob_out.pbData, blob_out.cbData);
        LocalFree(blob_out.pbData);
    }
}

// Linux keyring using libsecret
#include <libsecret/secret.h>

void store_key_in_keyring(const char *key_name, const uint8_t *key, size_t key_len) {
    SecretService *service = secret_service_get_sync(SECRET_SERVICE_NONE);
    
    SecretValue *value = secret_value_new((const gchar*)key, key_len, "text/plain");
    
    secret_service_store_sync(service, NULL, SECRET_SCHEMA_COMPAT_NETWORK,
                             NULL, key_name, NULL, NULL, SECRET_VALUE_ALL, value, NULL, NULL);
    
    g_object_unref(value);
    g_object_unref(service);
}

// Memory protection (mlock)
void protect_key_in_memory(uint8_t *key, size_t key_len) {
    // Lock memory to prevent swapping
    if (mlock(key, key_len) != 0) {
        perror("mlock failed");
        return;
    }
    
    // Mark memory as non-dumpable (Linux specific)
    if (madvise(key, key_len, MADV_DONTDUMP) != 0) {
        perror("madvise failed");
    }
    
    // Use secure memory allocation
    uint8_t *secure_key = (uint8_t*)malloc_secure(key_len);
    memcpy(secure_key, key, key_len);
    
    // Clear original key
    secure_zero_memory(key, key_len);
}
\end{lstlisting}

\subsection{Key Rotation and Revocation}

\subsubsection{Automated Key Rotation}

\begin{lstlisting}[style=cstyle, caption={Key Rotation System}]
typedef struct {
    uint8_t key_id[16];
    uint8_t key[32];
    time_t created;
    time_t expires;
    bool is_active;
} KeyEntry;

typedef struct {
    KeyEntry *keys;
    size_t capacity;
    size_t count;
    time_t rotation_interval;
} KeyManager;

// Initialize key manager
void key_manager_init(KeyManager *km, size_t capacity, time_t rotation_interval) {
    km->keys = calloc(capacity, sizeof(KeyEntry));
    km->capacity = capacity;
    km->count = 0;
    km->rotation_interval = rotation_interval;
}

// Generate and store new key
void key_manager_rotate(KeyManager *km) {
    time_t now = time(NULL);
    
    // Find expired keys
    for (size_t i = 0; i < km->count; i++) {
        if (km->keys[i].expires < now) {
            km->keys[i].is_active = false;
        }
    }
    
    // Generate new key
    if (km->count < km->capacity) {
        KeyEntry *new_key = &km->keys[km->count];
        
        // Generate unique key ID
        generate_uuid(new_key->key_id);
        
        // Generate cryptographic key
        generate_random_bytes(new_key->key, 32);
        
        new_key->created = now;
        new_key->expires = now + km->rotation_interval;
        new_key->is_active = true;
        
        km->count++;
        
        printf("Generated new key: ");
        for (int i = 0; i < 16; i++) printf("%02x", new_key->key_id[i]);
        printf(" (expires: %ld)\n", new_key->expires);
    }
}

// Get active key for encryption
int key_manager_get_active_key(KeyManager *km, uint8_t *key_id, uint8_t *key) {
    time_t now = time(NULL);
    
    for (size_t i = 0; i < km->count; i++) {
        if (km->keys[i].is_active && km->keys[i].expires > now) {
            memcpy(key_id, km->keys[i].key_id, 16);
            memcpy(key, km->keys[i].key, 32);
            return 0;
        }
    }
    
    return -1; // No active key found
}

// Background rotation thread
void *key_rotation_thread(void *arg) {
    KeyManager *km = (KeyManager*)arg;
    
    while (1) {
        sleep(km->rotation_interval);
        key_manager_rotate(km);
    }
    
    return NULL;
}
\end{lstlisting}

\subsubsection{Key Revocation}

\begin{lstlisting}[style=cstyle, caption={Key Revocation System}]
typedef struct {
    uint8_t key_id[16];
    time_t revoked_at;
    char reason[256];
} RevokedKey;

typedef struct {
    RevokedKey *revoked_keys;
    size_t count;
    size_t capacity;
} RevocationList;

// Revoke key
void revoke_key(RevocationList *rl, const uint8_t *key_id, const char *reason) {
    if (rl->count < rl->capacity) {
        RevokedKey *rk = &rl->revoked_keys[rl->count];
        
        memcpy(rk->key_id, key_id, 16);
        rk->revoked_at = time(NULL);
        strncpy(rk->reason, reason, 255);
        rk->reason[255] = '\0';
        
        rl->count++;
        
        printf("Key revoked: ");
        for (int i = 0; i < 16; i++) printf("%02x", key_id[i]);
        printf(" (%s)\n", reason);
    }
}

// Check if key is revoked
bool is_key_revoked(RevocationList *rl, const uint8_t *key_id) {
    for (size_t i = 0; i < rl->count; i++) {
        if (memcmp(rl->revoked_keys[i].key_id, key_id, 16) == 0) {
            return true;
        }
    }
    return false;
}

// Publish revocation list (CRL)
void publish_crl(RevocationList *rl, const char *filename) {
    FILE *f = fopen(filename, "w");
    if (!f) return;
    
    fprintf(f, "# Certificate Revocation List\n");
    fprintf(f, "# Generated: %s", ctime(&(time_t){time(NULL)}));
    
    for (size_t i = 0; i < rl->count; i++) {
        RevokedKey *rk = &rl->revoked_keys[i];
        fprintf(f, "REVOKED: ");
        for (int j = 0; j < 16; j++) fprintf(f, "%02x", rk->key_id[j]);
        fprintf(f, " %ld %s\n", rk->revoked_at, rk->reason);
    }
    
    fclose(f);
}
\end{lstlisting}

\subsection{Key Escrow and Recovery}

\subsubsection{Secret Sharing}

Shamir's Secret Sharing membagi kunci menjadi beberapa bagian yang dapat direkonstruksi dengan threshold tertentu \cite{shamir1979}.

\begin{lstlisting}[style=cstyle, caption={Shamir's Secret Sharing}]
#include <stdint.h>
#include <stdlib.h>

#define FIELD_SIZE 257  // Prime field

// Polynomial evaluation
uint16_t evaluate_polynomial(uint16_t *coeffs, int degree, uint16_t x) {
    uint16_t result = 0;
    
    for (int i = degree; i >= 0; i--) {
        result = (result * x + coeffs[i]) % FIELD_SIZE;
    }
    
    return result;
}

// Generate shares
void generate_shares(uint16_t secret, int threshold, int num_shares, 
                    uint16_t *x_coords, uint16_t *y_coords) {
    // Generate random polynomial
    uint16_t coeffs[threshold];
    coeffs[0] = secret;
    
    for (int i = 1; i < threshold; i++) {
        coeffs[i] = random_uint16() % FIELD_SIZE;
    }
    
    // Evaluate polynomial at different x coordinates
    for (int i = 0; i < num_shares; i++) {
        x_coords[i] = i + 1;  // Use 1, 2, 3, ... as x coordinates
        y_coords[i] = evaluate_polynomial(coeffs, threshold - 1, x_coords[i]);
    }
}

// Lagrange interpolation for reconstruction
uint16_t reconstruct_secret(uint16_t *x_coords, uint16_t *y_coords, 
                           int threshold) {
    uint16_t secret = 0;
    
    for (int i = 0; i < threshold; i++) {
        uint16_t numerator = 1;
        uint16_t denominator = 1;
        
        // Calculate Lagrange basis polynomial
        for (int j = 0; j < threshold; j++) {
            if (i != j) {
                numerator = (numerator * (-x_coords[j])) % FIELD_SIZE;
                denominator = (denominator * (x_coords[i] - x_coords[j])) % FIELD_SIZE;
            }
        }
        
        // Add contribution
        uint16_t contribution = y_coords[i] * numerator * mod_inverse(denominator, FIELD_SIZE);
        secret = (secret + contribution) % FIELD_SIZE;
    }
    
    return secret;
}

// Usage example
void secret_sharing_example() {
    uint16_t secret = 12345;  // Secret to protect
    int threshold = 3;       // Need 3 shares to reconstruct
    int num_shares = 5;       // Generate 5 shares total
    
    uint16_t x_coords[5], y_coords[5];
    
    // Generate shares
    generate_shares(secret, threshold, num_shares, x_coords, y_coords);
    
    printf("Generated shares:\n");
    for (int i = 0; i < num_shares; i++) {
        printf("Share %d: (%d, %d)\n", i+1, x_coords[i], y_coords[i]);
    }
    
    // Reconstruct with first 3 shares
    uint16_t reconstructed = reconstruct_secret(x_coords, y_coords, threshold);
    printf("Reconstructed secret: %d\n", reconstructed);
}
\end{lstlisting}

\subsection{Best Practices and Guidelines}

\textbf{Key Management Best Practices}:
\begin{itemize}
  \item \textbf{Use CSPRNG}: Generate keys with cryptographically secure random numbers
  \item \textbf{Minimum Key Lengths}: AES-256, RSA-4096, ECC-384 for new systems
  \item \textbf{Regular Rotation}: Rotate keys every 90-365 days for high-value systems
  \item \textbf{Secure Storage}: Use HSM or encrypted storage for private keys
  \item \textbf{Access Control}: Implement principle of least privilege
  \item \textbf{Audit Logging}: Log all key operations
  \item \textbf{Backup}: Secure backup of key material
  \item \textbf{Revocation}: Have clear revocation procedures
\end{itemize}

\begin{catatan}
Key management is often the weakest link in cryptographic systems. Even the strongest algorithms can be compromised through poor key management practices. Implement defense in depth: use hardware security modules, regular audits, and proper procedures.
\end{catatan}
